% =========================
% Cas pratique 1
% =========================

\section[Cas 1]{Cas pratique 1 : Les constats d’huissier}

\begin{frame}{Question 1}

\begin{block}{}
L'art. \textbf{L615-5} al. 1 CPI : "tous les moyens".
\end{block}

\begin{table}[h]
\centering
\begin{tabular}{l l l}
    Quand? & Moyen de preuve \\
    \hline
    av. action en cf. & constat (- force probante: témoignage) \\
    au début d'une action en cf. & saisie cf. \\
    à la fin d'une action en cf. & droit de l'information \\
\end{tabular}
\end{table}
\end{frame}

\begin{frame}{Question 2}

\begin{block}{}
L'art. \textbf{D631-2} CPI : "Tribunal Judiciaire de Paris".
\\ Principe de droit: specialia generalibus derogant.
\end{block}

\begin{itemize}
    \item En l'espèce, Tribnual Jurdicaire de Saint-Dié-des-Vosges.
    De plus, L615-5 (spécifique) déroge à Art. 145 CPC (général).
    \item En conclusion, mauvaise compétence et fondement nul. 
    Donc, ordonnance nulle.
    Par conséquent, PV 14/11/2023 nul.
\end{itemize}
\end{frame}

\begin{frame}{Question 3}

\begin{itemize}
    \item En l'espèce, \begin{itemize}
        \item description et capture pages web pas ds pt. 1 ordonnance
        \item 1 PV ("je reprends"), pas plusieurs dans pt. 2 ordonnance.
        \item stagiaire, pas homme de l'art, pt. 3 ordonnance.
    \end{itemize}
    \item En conclusion, au moins pour ces raisons,
    le PV 15/11/2023 est nul.
\end{itemize}

\end{frame}

\begin{frame}{Question 4}

\begin{itemize}
    \item En l'espèce, analyse pas dans ordonnance.
    \item En conclusion, au moins pour cette raison,
    le PV 17/11/2023 est nul.
\end{itemize}

\end{frame}